\needspace{9\baselineskip}
\begin{center}
\pdfbookmark[1]{5. Zu Art. 78 II.}{zusatz-005}
[5. Zu Art. 78 II.]
\end{center}
Directe Auflösung der Gleichungen
\[
\begin{aligned}
1)\quad&\sin(P-A)=k\sin(Q-B)\\
2)\quad&\sin(P-A')=k'\sin(Q-B')
\end{aligned}
\]
[wo \(A,A',B,B',k,k'\) gegebene Grössen und \(P,Q\) zu bestimmen sind].

Es folgt daraus
\[
\begin{aligned}
\operatorname{tang}\tfrac12(P+Q-A-B)&=\frac{k+1}{k-1}\operatorname{tang}\tfrac12(P-Q-A+B)\\
\operatorname{tang}\tfrac12(P+Q-A'-B')&=\frac{k'+1}{k'-1}\operatorname{tang}\tfrac12(P-Q-A'+B')
\end{aligned}
\]
oder wenn wir
\[
\begin{aligned}
\tfrac12(P+Q-A-B)&=x, &\tfrac12(P-Q-A+B)&=y\\
\tfrac12(A+B-A'-B')&=\alpha, &\tfrac12(A-B-A'+B')&=\beta\\
\frac{k+1}{k-1}&=m, &\frac{k'+1}{k'-1}&=n
\end{aligned}
\]
setzen:
\[
\begin{aligned}
3)\quad&\operatorname{tang}x=m\operatorname{tang}y\\
4)\quad&\operatorname{tang}(x+\alpha)=n\operatorname{tang}(y+\beta)
\end{aligned}
\]
\SourceLeaf{304}
oder in imaginärer Form
\[
\frac{e^{ix}-e^{-ix}}{e^{ix}+e^{-ix}}=m\frac{e^{iy}-e^{-iy}}{e^{iy}+e^{-iy}}
\]
[und analog für Gleichung 4] oder
\[
e^{2ix}=\frac{(1+m)e^{iy}+(1-m)e^{-iy}}{(1-m)e^{iy}+(1+m)e^{-iy}}
\]
und
\[
e^{2i(x+\alpha)}=\frac{(1+n)e^{i(y+\beta)}+(1-n)e^{-i(y+\beta)}}{(1-n)e^{i(y+\beta)}+(1+n)e^{-i(y+\beta)}}\mathpunct{\raisebox{.45ex}{;}}
\]
folglich
\[
\begin{aligned}
&\bigl((1+m)e^{iy}+(1-m)e^{-iy}\bigr)\bigl((1-n)e^{i(y+\beta)}+(1+n)e^{-i(y+\beta)}\bigr)e^{i\alpha}\\
={}&\bigl((1-m)e^{iy}+(1+m)e^{-iy}\bigr)\bigl((1+n)e^{i(y+\beta)}+(1-n)e^{-i(y+\beta)}\bigr)e^{-i\alpha}
\end{aligned}
\]
oder wenn man entwickelt und mit 2 dividirt:
\[
\begin{aligned}
&e^{i(2y+\beta)}\bigl((m-n)\cos\alpha+i(1-mn)\sin\alpha\bigr)\\
&\quad-e^{-i(2y+\beta)}\bigl((m-n)\cos\alpha-i(1-mn)\sin\alpha\bigr)\\
={}&2i\bigl\{(m+n)\cos\alpha\sin\beta-(1+mn)\sin\alpha\cos\beta\bigr\}.
\end{aligned}
\]

Führt man also die Hülfsgrössen \(g,G,h,H\) ein, so dass
\[
\begin{aligned}
(m-n)\cos\alpha&=g\cos G, &(m+n)\cos\alpha&=h\cos H\\
(1-mn)\sin\alpha&=g\sin G, &(1+mn)\sin\alpha&=h\sin H,
\end{aligned}
\]
so wird diese Gleichung
\[
5)\quad g\sin(2y+\beta+G)=h\sin(\beta-H),
\]
woraus die beiden Werthe von \(y\) sich ergeben. Die correspondirenden Werthe von \(x\) finden sich dann aus (3) oder (4).

Man kann auch setzen
\[
\begin{aligned}
(k'-k)\cos\alpha&=g^*\cos G, &(kk'-1)\cos\alpha&=h^*\cos H\\
-(k'+k)\sin\alpha&=g^*\sin G, &(kk'+1)\sin\alpha&=h^*\sin H
\end{aligned}
\]
und hat dann
\[
g^*\sin(2y+\beta+G)=h^*\sin(\beta-H).
\]
[Übrigens ist hier
\[
\begin{aligned}
g^*&=\tfrac12(k-1)(k'-1)g\\
h^*&=\tfrac12(k-1)(k'-1)h.]
\end{aligned}
\]
\begin{center}\rule{.18\textwidth}{.35pt}\end{center}
